\chapter{Аудит фермионной cross-билинейности и нечётной статистики}
\label{ch:v10-m4-fermionic-cross-odd-audit}

Условный determinant-механизм требует одновременно cross-оператор
правильного $M_4$-типа, нечётную статистику, отрицательную восприимчивость,
gauge/KMS-совместимость и независимо выбранный coupling. Проверены
двенадцать источников; матрица шести критериев имеет ранг $6$, но
\begin{equation}
 N_{\rm pass}=0/12,\qquad N_{\rm max}=5/6.
\end{equation}

KO6 particle/conjugate-пара наследуется и фермионна, однако для
$Q=\operatorname{diag}(I_2,-I_2)$ и cross-инволюции $X$ выполняется
\begin{equation}
 \rank[Q,X]=4.
\end{equation}
Значит, этот flip смешивает противоположные заряды. Reduced Pfaffian
различает две ориентации знаками $(-1,+1)$, но полный KO6 pair даёт
$(+1,+1)$ и стирает относительный знак.

Наследуемая секторная grading имеет odd-rank $2$, тогда как determinant
требует all-odd rank $4$; дефект имеет ранг $2$. Формальный parity shift
исправляет счёт, но не выводится из физической секторной градуировки.

Единственный gauge-совместимый почти проходящий маршрут --- равнозарядная
каллиасова twist-двойка: её charge-коммутатор равен нулю, и она условно
поддерживает cross-билинейность, нечётность и determinant-знак. Однако
унаследованное отображение каллиасова носителя размерности $60$ в
$M_4$-модуль равно нулю и имеет ранг $0$.

\begin{theorem}[Тройной разрыв fermionic cross origin]
KO6 route наследуется, но нарушает заряд; reduced Pfaffian теряет знак при
полном Real-удвоении; равнозарядный Callias route совместим, но не имеет
унаследованного вложения в $M_4$. Поэтому ни один из двенадцати кандидатов
не закрывает determinant-origin.
\end{theorem}

\begin{nogocriterion}[Operator grading не равна Grassmann parity]
Секторная градуировка ранга $2+2$ не назначает все четыре координаты
нечётными. Ручное объявление all-odd parity или cross-Grassmann pair
является новым statistics datum, пока не предъявлен общий типизированный
носитель или BRST/BV-родитель.
\end{nogocriterion}

\begin{protocol}\begin{enumerate}
\item проверено кандидатов: \textbf{$12$};
\item ранг критериев: \textbf{$6/6$};
\item полных проходов: \textbf{$0/12$};
\item KO6 charge defect: \textbf{rank $4$};
\item inherited/required odd rank: \textbf{$2/4$};
\item Callias--$M_4$ embedding: \textbf{rank $0$};
\item следующий гейт: \textbf{типизированное каллиасово равнозарядное вложение}.
\end{enumerate}\end{protocol}

Машинный сертификат сохранён в
\path{s2t_v10_cell_birth_four_volume_induced_newton_breathing_anomaly_hopf_reservoir_m4_fermionic_determinant_cross_bilinear_odd_statistics_candidate_audit_gate_results.json}.